% HugMesh Games — Architecture Reference
%
% Build:
%   make -C docs/publishing pdf
%   ./docs/publishing/scripts/build-pdf.sh
%
% Output: docs/publishing/architecture/build/architecture.pdf
% Cleaned + ignored — see docs/publishing/.gitignore.

\documentclass[11pt,a4paper,openany]{book}

\usepackage{hugmesh-theme}
\usepackage{microtype}
\usepackage[hidelinks]{hyperref}
\usepackage{tabularx}
\usepackage{longtable}
\usepackage{booktabs}
\usepackage{array}
\usepackage{datetime}
% pdfpages lets us drop the auto-sized DB-diagram PDFs onto their own
% over-sized pages so a reviewer can pinch-zoom in any PDF viewer without
% the included content getting rasterised down to thumbnail size.
\usepackage{pdfpages}
% afterpage queues multiple per-page hooks (unlike \AddToHookNext{shipout}
% which only retains the most-recent call); used by the schema-diagram
% includes to set \pdfpagewidth / \pdfpageheight independently per
% diagram page.
\usepackage{afterpage}

\hypersetup{
  pdftitle={HugMesh Games — Architecture Reference},
  pdfauthor={HugMesh Systems},
  pdfsubject={Architecture, build pipeline, and API reference for the
              HugMesh Games portal (hugmesh.games)},
  pdfkeywords={Angular, Spring Boot, Godot, WebAssembly, Module Federation,
               GraalVM, Casdoor, JWT, R2DBC, PostgreSQL, Cassandra,
               docker-compose, HugMesh Games},
}

\begin{document}

% \sloppy + \emergencystretch (declared in hugmesh-theme.sty) let
% LaTeX stretch interword spaces a bit more than the strict default so
% paragraphs that contain long unbreakable tokens (URLs, file paths,
% camelCase identifiers) lay out without overflowing the right margin.
\sloppy

% ── Cover page ────────────────────────────────────────────────────────────
\begin{titlepage}
\thispagestyle{empty}
\noindent
\begin{tikzpicture}[remember picture,overlay]
  \fill[snapBg] (current page.south west) rectangle (current page.north east);
  \node[anchor=north west,inner sep=0pt] at ($(current page.north west)+(0,-0.0\paperheight)$) {
    \begin{tikzpicture}
      \fill[snapAccent,opacity=0.18]
        (0,0) circle (8cm);
      \fill[snapAccentTwo,opacity=0.10]
        (4,-2) circle (5cm);
    \end{tikzpicture}
  };
  \node[anchor=south east,inner sep=0pt] at ($(current page.south east)+(0,0.05\paperheight)$) {
    \begin{tikzpicture}
      \fill[snapAccent,opacity=0.12]
        (0,0) circle (6cm);
    \end{tikzpicture}
  };
\end{tikzpicture}

\vspace*{4cm}
\noindent
{\sffamily\bfseries
  {\color{snapAccentTwo}\fontsize{18}{22}\selectfont HUGMESH GAMES}\\[3mm]
  {\color{white}\fontsize{42}{50}\selectfont Architecture}\\[2mm]
  {\color{white}\fontsize{42}{50}\selectfont Reference}\par}

\vspace{1.5cm}
\noindent\color{white!75!snapAccentTwo}\rule{0.4\linewidth}{1.2pt}

\vspace{0.8cm}
\noindent
{\sffamily\color{white!85}\large
The federated Angular host, the Spring Boot BFF compiled to a GraalVM
native image, the per-game Godot WebAssembly bundles, and the hot-upload
pipeline that links them.}

\vfill
\noindent
{\sffamily\color{white!70}\small
\snapprojectname{} portal --- hugmesh.games\\
Generated: \snapbuilddate{} \snapbuildtime}

\end{titlepage}

% ── Front matter ──────────────────────────────────────────────────────────
\frontmatter
\pagestyle{empty}

\noindent
\begin{snapnote}[title={About this document}]
This is a typeset companion to the canonical
\href{https://github.com/hugmesh-net/hugmesh-games}{repository
\texttt{README.md}} and the
\href{https://github.com/hugmesh-net/hugmesh-games/blob/main/docs/PLAN-module-federation-and-native.md}{post-federation
\texttt{PLAN}} document. Where it disagrees with the source tree, the source
tree wins — every diagram, table, and snippet here is generated from a
specific commit and may drift between rebuilds.

To rebuild it from a fresh checkout:
\begin{lstlisting}[language=bash]
make docs-pdf       # or: make -C docs/publishing pdf
\end{lstlisting}
The build script (\reposrc{docs/publishing/scripts/build-pdf.sh}) runs
\texttt{pdflatex} twice to resolve cross-references and the table of
contents. The build is idempotent and produces no committed artefacts —
all output lives in \texttt{docs/publishing/architecture/build/} which is
git-ignored.
\end{snapnote}

\vspace{1em}
\tableofcontents
\clearpage

% ── Main matter ───────────────────────────────────────────────────────────
\mainmatter
\pagestyle{fancy}

\chapter{Overview}

\textbf{HugMesh Games} is a browser-native multi-game portal: a single
Angular shell loads multiple Godot 4 WebAssembly titles on demand, with a
Spring Boot BFF streaming game state and controls over Server-Sent Events.
Authentication runs through Casdoor (OAuth2/OIDC) behind a federated login
remote.

\section{Three sentences}

The host shell is a \textbf{native-federation} Angular host that loads its
\texttt{login}, \texttt{admin}, and per-game routes from sibling bundles
served by per-remote nginx containers on the compose network. The BFF
compiles to either a fat JAR or a \textbf{GraalVM native image}, exposes
the SSE stream, the federation manifest, the upload pipeline, and the
per-game launch / access endpoints, and persists state in PostgreSQL with
a Cassandra cross-pod cache. Each Godot game lives in its own project
under \snapfile{games/<slug>/}, exports a self-contained WASM bundle, and
is delivered to the browser either through the BFF asset path or through
the dedicated \texttt{wasm-fragment} nginx remote with the
COOP/COEP/CORP headers WASM games need.

\section{What's in the box}

\begin{itemize}
  \item \textbf{Angular host} (\reposrc{frontend/}) — federation host plus
        four sibling remotes (\texttt{login}, \texttt{admin},
        \texttt{games-three-wizards}, \texttt{games-dodecahedron}) and
        the asset-tier \texttt{wasm-fragment}.
  \item \textbf{Spring Boot BFF} (\reposrc{backend/}) — WebFlux + R2DBC,
        Liquibase via JDBC, Casdoor JWT verification, GraalVM-ready.
  \item \textbf{Per-game Godot projects} (\reposrc{games/}) — currently
        \texttt{three-wizards-in-a-trenchcoat} (cave + field) and
        \texttt{dodecahedron-demo} (chrome-less single-scene tech demo).
  \item \textbf{Reusable Rust GDExtension} (\reposrc{extensions/hugmesh-bridge/})
        — native \texttt{.so} for the editor, \texttt{wasm32-unknown-emscripten}
        for the web build, with a typed signal bus that any new game can
        consume.
  \item \textbf{wasm-fragment nginx} (\reposrc{wasm-fragment/}) — bakes
        \snapfile{backend/src/main/resources/assets/levels/} and serves
        Godot bundles at \texttt{/\_remotes/wasm-fragment/<slug>/index.html}
        with COOP \texttt{same-origin}, COEP \texttt{require-corp}, and
        CORP \texttt{cross-origin}.
  \item \textbf{Casdoor} (\reposrc{casdoor/}) — never exposed to the host;
        the BFF reverse-proxies the login flow so the admin UI stays
        internal.
  \item \textbf{PostgreSQL + Cassandra} — Postgres holds the canonical
        relational state; Cassandra is the cross-pod movement / session
        cache and a write-through audit log.
\end{itemize}

\section{The plan that produced this shape}

The post-federation topology landed under
\reposrc{docs/PLAN-module-federation-and-native.md}, an eight-section plan
spanning module federation (A), Godot editor plugin (B), hot upload (C),
GraalVM native-image (D), JWT roles + per-game access (E), DB-backed
federation discovery (F), per-game project layout (G), and cryptographic
launch UIDs (H). Every section is at a deliberate end state at the head
of the \reposrc{feature/module-federation-and-native} branch; this
document is the printed companion that explains why each piece exists and
how they fit.

\begin{snapnote}[title={Audience}]
This reference is for engineers who need to extend, deploy, or debug the
portal. The administrator manual, the entry-level tour, and the Godot
developer walkthrough are separate slide decks under
\reposrc{docs/publishing/slides/} and target their respective audiences
without overlap.
\end{snapnote}

\section{What's new since the post-federation snapshot}

Several feature batches have landed on top of the eight-section
post-federation plan; each is covered in detail in its own chapter:

\begin{itemize}
  \item \textbf{Multiplayer sessions.} Per-session SSE sinks
        (\snapfile{ConcurrentHashMap<String, SessionSinks>}) layered over
        the legacy global sink; \texttt{applyControl(event, characterId,
        sessionId)} is the multiplayer-aware overload.
        \snapfile{SessionIsolationTest} pins the cross-session
        non-leakage contract. See chapter~\ref{ch:runtime}.
  \item \textbf{Friends, blocks, session visibility.} Single
        \texttt{user\_relationships} graph table backs both mutual
        friendship and one-directional blocks; \texttt{game\_sessions
        .visibility} is one of \texttt{public} / \texttt{friends} /
        \texttt{private}. The lobby filter runs against the caller's
        friend graph + blocklist. Decision 015 has the full design.
  \item \textbf{Mod and publisher roles with per-game scope.} Casdoor
        owns role membership; Postgres owns the per-game scope in
        \texttt{mod\_game\_scopes} and \texttt{publisher\_game\_scopes}.
        See chapter~\ref{ch:datamodel}.
  \item \textbf{RPG + VN makers.} Two reactive Spring Boot side services
        (\reposrc{rpg-builder/} and \reposrc{vn-builder/}) with their
        own Postgres instances, federated frontend remotes at
        \texttt{/rpg-maker} and \texttt{/vn-maker}, and a build pipeline
        that POSTs the produced bundle to the BFF's
        \texttt{/api/admin/games/upload} endpoint with the publisher's
        bearer. Decision 018 covers the side-service architecture.
  \item \textbf{Production rate-limit zones.} Six nginx
        \texttt{limit\_req\_zone} declarations in
        \snapfile{frontend/nginx.conf} (and the TLS twin) sized for
        $\sim$10 000 concurrent users. The BFF still runs its per-pod
        \snapfile{RateLimitFilter} + Cassandra cross-pod aggregator
        behind the edge. Decision 019.
  \item \textbf{LAN access + TLS overlay.} \snapfile{lan.properties}
        env-file with literal \texttt{FRONTEND\_BIND=0.0.0.0:} values;
        \texttt{make lan-up\{,-tls\}} wraps the stack with self-signed
        TLS (cert SAN includes detected LAN IPs).
  \item \textbf{Crypto.mintUid() + Liquibase 036.} 43-char URL-safe
        base64 SHA3-256 minter replacing \texttt{UUID.randomUUID()} for
        every UI-facing UID; Liquibase 036 widens session-id columns to
        \texttt{varchar(64)} to fit the new shape without disturbing
        pre-existing UUID rows.
\end{itemize}

\chapter{Runtime architecture}

The runtime data flow has three tiers. Tier one is the player's browser:
the federated Angular host, the active remote module, and the WASM iframe.
Tier two is the per-remote nginx fleet: the host, four route-or-game
remotes, and the asset-tier \texttt{wasm-fragment}. Tier three is the BFF
plus its persistence: Spring Boot WebFlux on top of Postgres and Cassandra,
behind Casdoor for authentication.

\section{Browser-side flow}

\begin{figure}[!h]
\centering
\snapdiagram{diagrams/runtime-flow.tex}
\caption{Runtime data flow from a click on the home page to a frame
rendered inside the Godot iframe. \texttt{SCENE\_TYPE\_TO\_GAME\_SLUG}
forks the renderer; federated games mint a \texttt{launch\_uid} for
operational visibility; \texttt{TW\_LEVEL\_TRANSITION} \texttt{postMessage}
is caught at the host level so it works regardless of which renderer is
mounted.}
\end{figure}

The host shell's \snapfile{main.ts} is intentionally split into two
modules. \snapfile{main.ts} runs the
\textbf{native-federation \texttt{initFederation(...)}} bootstrap with the
manifest fetched from \texttt{GET /api/federation/manifest}, then stashes
\texttt{loadRemoteModule} on \texttt{globalThis.\_\_SNAPDRAGON\_LOAD\_REMOTE\_MODULE\_\_}
and only then dynamically imports \snapfile{bootstrap.ts}, which calls
\texttt{bootstrapApplication} with the host's \texttt{AppComponent}. This
ordering is non-negotiable: any \texttt{@angular/*} static import resolved
before \texttt{initFederation} runs will fail because the import map
that's about to be installed isn't there yet.

\begin{snapnote}[title={Why \texttt{globalThis} and not a service?}]
The federation runtime can't be reached through Angular's DI without
making \texttt{@angular-architects/native-federation} a static import
into a route-handler chunk, which puts it on karma's module graph and
breaks tests that don't load the federation runtime at all. Stashing the
function on \texttt{globalThis} after bootstrap keeps the package out of
karma's static graph and out of every route bundle that doesn't need it.
\end{snapnote}

\section{Three-tier topology}

\begin{figure}[!h]
\centering
\snapdiagram{diagrams/topology.tex}
\caption{Compose-network view. Each box is a Docker container; arrows are
HTTP routes proxied by the host nginx at \snapfile{frontend/nginx.conf}.
The BFF is the only container that talks to Postgres / Cassandra /
Casdoor.}
\end{figure}

\subsection{The host shell}

The host nginx serves the host bundle plus a set of \texttt{location
/\_remotes/<slug>/} blocks that proxy each remote's
\texttt{remoteEntry.json} and chunks. For \texttt{wasm-fragment} the host
strips upstream COOP/COEP duplicates with \texttt{proxy\_hide\_header} and
re-emits the canonical set so the WASM iframe loads with isolated
headers.

\subsection{The remotes}

\begin{tabularx}{\linewidth}{l l X}
\toprule
\textbf{Slug} & \textbf{Kind} & \textbf{What it exposes / does} \\
\midrule
\texttt{login} & \texttt{remote-shell} & \texttt{./LoginComponent} —
mounted via \texttt{loadComponent} on the \texttt{/login} route. Inline
\texttt{POST /api/auth/token} + \texttt{sessionStorage.tw\_auth} write,
mirrors the host \texttt{AuthService.StoredAuth} shape. \\
\texttt{admin} & \texttt{remote-shell} & \texttt{./AdminShell} +
\texttt{./UploadShell} — both expose under
\texttt{federation\_modules.slug = 'admin'} so \texttt{/admin} and
\texttt{/admin/games} share a build. \\
\texttt{games-three-wizards} & \texttt{remote-game} &
\texttt{./GameRoot} — Web Component (\texttt{createCustomElement}) with
\texttt{@Input levelSlug + levelId + launchUid}. Renders cave and field. \\
\texttt{games-dodecahedron} & \texttt{remote-game} &
\texttt{./GameRoot} — same pattern, single chrome-less scene. \\
\texttt{wasm-fragment} & \texttt{fragment} & Asset-tier nginx; bakes
\snapfile{backend/src/main/resources/assets/levels/} and serves at
\texttt{/\_remotes/wasm-fragment/<slug>/index.html}. \\
\bottomrule
\end{tabularx}

\subsection{The BFF}

\snapfile{backend/} runs as a single Spring Boot 3.4 application with both
WebFlux (R2DBC, runtime path) and JDBC (Liquibase, blocking) wired in.
Spring auto-config skips \texttt{DataSourceAutoConfiguration} when an
R2DBC \texttt{ConnectionFactory} is present, so
\snapfile{JdbcDataSourceConfig} declares an explicit
\texttt{@Primary HikariDataSource} of pool size 3 just for Liquibase to
acquire on boot.

\section{Authentication and access}

Casdoor is reachable only from inside the compose network. The host BFF
proxies the login flow at \texttt{/api/auth/*} and verifies JWTs against
the anonymous \texttt{/.well-known/jwks} endpoint (not
\texttt{/api/get-jwks}, which requires admin auth and would 401 on every
request). Per-game visibility is enforced by the \texttt{game\_access}
table — see chapter~\ref{ch:datamodel}.

\section{Server-sent events}

\snapfile{GameStreamController} merges \emph{eight} SSE channels onto a
single \texttt{Flux<ServerSentEvent>}: seven \texttt{Sinks.Many} from
\snapfile{GameStateService} (\texttt{movement}, \texttt{world},
\texttt{character}, \texttt{scene}, \texttt{input}, \texttt{audio},
\texttt{wasm-console}) plus one from \snapfile{GameTimeService}
(\texttt{time}), with a 15-second \texttt{Flux.interval} keepalive
that fires immediately on subscribe so the response headers flush at
connect time and the browser \texttt{EventSource.onopen} fires in
milliseconds.

The latter three channels are part of the generalised input surface
introduced for gamepad / mouse / keyboard / touch / audio support
(see chapter~\ref{ch:datamodel}'s reference-table description and
chapter~\ref{ch:api}'s \emph{Input events} + \emph{Diagnostics} tags).
Subscribers opt into individual channels via the EventSource
event-listener model (\texttt{source.addEventListener('input', cb)});
merging onto the single connection keeps the host page from needing
N EventSources to see everything.

\begin{snapcaution}[title={The \texttt{autoCancel} trap}]
\texttt{Sinks.many().multicast().onBackpressureBuffer(N)} defaults to
\texttt{autoCancel=true}. When every subscriber disconnects (iframe
remount, EventSource reconnect, network blip), the sink terminates
\textbf{permanently} and all subsequent \texttt{tryEmitNext} calls
silently drop. Every sink in this project passes \texttt{autoCancel=false}
explicitly. If a freshly-deployed sink \enquote{works once and then
silently dies}, this is the cause.
\end{snapcaution}

\section{Per-session SSE sink isolation}

Multiplayer sessions add a per-session sink layer on top of the
legacy global sinks. \snapfile{GameStateService} maintains a
\texttt{ConcurrentHashMap<String, SessionSinks>} lazy-allocated by
\texttt{sessionId}; when a control event arrives with a session id,
\texttt{applyControl(event, characterId, sessionId)} fans out to the
matching per-session bundle \emph{and} the global sink. A subscriber
to \texttt{GET /api/game/stream?sessionId=$\alpha$} sees only session
$\alpha$'s events; the admin firehose at
\texttt{/api/admin/stream} keeps the cross-session view it had before.

\begin{figure}[!h]
\centering
\snapdiagram{diagrams/sessions-topology.tex}
\caption{Per-session SSE sink isolation. Events emitted on session
$\alpha$'s sinks NEVER reach session $\beta$'s subscribers and vice
versa; the global sink keeps the admin firehose intact. The
\reposrc{backend/src/test/java/games/snapdragon/bff/service/SessionIsolationTest.java}
integration test pins the dashed-arrow non-leakage in both directions.}
\end{figure}

The new \emph{view} channel carries per-player camera + orientation
samples ($\sim$10\,Hz) emitted by the Rust GDExtension's
\texttt{view\_egress} signal. Admins watch via
\texttt{GET /api/admin/sessions/\{id\}/stream} which merges movement,
world, character, and view onto a single per-session connection.
Without per-session keying, two concurrent sessions on the same pod
would cross-contaminate every fan-out — so the keyed map is the
load-bearing piece.

\section{Maker side services}

The RPG and VN makers each ship as their own reactive Spring Boot
service alongside the BFF, with their own Postgres tenants:

\begin{itemize}
  \item \snapfile{rpg-builder/} --- Mudlet-shaped editor data model
        (projects, areas, rooms, exits, items, NPCs, scripts) backed
        by the \texttt{postgres-rpg-builder} container. REST surface
        at \texttt{/api/rpg/*}.
  \item \snapfile{vn-builder/} --- godot-2d-visual-novel-shaped data
        model (projects, acts, characters, backgrounds, scenes, lines,
        choices) backed by the \texttt{postgres-vn-builder} container.
        REST surface at \texttt{/api/vn/*}.
\end{itemize}

Each side service's build worker walks the project tree, emits a Godot
4 source tree with \snapfile{snapdragon\_bridge.gdextension} baked in
(multiplayer-aware out of the box), runs the Godot CLI per area or act,
zips the result, and POSTs to the BFF's
\texttt{/api/admin/games/upload} endpoint \emph{with the originating
publisher's bearer token}. The BFF's existing
\texttt{publisher\_game\_scopes} check enforces who can ship which
slug, so a single compromised maker can't spam the catalog. Decision
018 covers the side-service architecture; the per-database schemas
live in the \emph{Schema diagrams} appendix at the end of this
document.

\section{Mod / publisher per-game scope}

Authorisation for mod- and publisher-only endpoints reads two layers,
both required:

\begin{itemize}
  \item Casdoor membership in \texttt{ROLE\_MOD} or
        \texttt{ROLE\_PUBLISHER} (claim is in the JWT).
  \item One row in \texttt{mod\_game\_scopes} /
        \texttt{publisher\_game\_scopes} per (casdoorId, gameSlug)
        pair (Postgres).
\end{itemize}

Admins manage both layers via \texttt{/api/admin/role-grants/*};
publishers can grant mod on \emph{games they own} via
\texttt{/api/publisher/grant-mod}, which walks the same write path but
enforces caller-owns-target-game before mutating. The full rejection
matrix is pinned by
\reposrc{e2e/features/security/mod-grant-gates.feature} on the BDD side
and a \texttt{@WebFluxTest} on the controller side.

\chapter{Module Federation}

The host shell loads \emph{everything} above the WASM iframe through
\texttt{@angular-architects/native-federation}. Adding a new
route-bearing remote means adding a row to \texttt{federation\_modules},
shipping a per-remote bundle, and pointing \texttt{loadComponent} at the
exposed module. Adding a new game means the same plus an
\texttt{<app-game>-remote>} Web Component the host can place in its
template.

\section{The two ordering rules}

\begin{enumerate}
  \item \textbf{\texttt{initFederation} before \texttt{bootstrapApplication}.}
        \snapfile{frontend/src/main.ts} dynamic-imports
        \snapfile{bootstrap.ts} after the federation runtime resolves.
        Static imports of \texttt{@angular/*} from \texttt{main.ts} would
        be evaluated before the import map is installed and would
        therefore fail to resolve.
  \item \textbf{Remote-loader access through \texttt{globalThis}.}
        Route handlers reach the federation runtime through
        \texttt{globalThis.\_\_SNAPDRAGON\_LOAD\_REMOTE\_MODULE\_\_}, set
        by \texttt{main.ts} once \texttt{initFederation} resolves. This
        keeps \texttt{@angular-architects/native-federation} out of
        karma's static graph and out of every route-handler bundle.
\end{enumerate}

\section{The manifest}

\begin{figure}[!h]
\centering
\snapdiagram{diagrams/federation-bootstrap.tex}
\caption{Federation bootstrap. The host fetches the manifest from the BFF
\emph{before} \texttt{bootstrapApplication}, then constructs the import
map, then mounts the host \texttt{AppComponent}. New remotes appear by
inserting rows into \texttt{federation\_modules} — no host redeploy
required.}
\end{figure}

The manifest endpoint is
\texttt{GET /api/federation/manifest} — DB-backed, public on purpose
because the host needs the \texttt{login} remote URL before any user has
authenticated. It returns the active rows of \texttt{federation\_modules}
ordered by \texttt{kind, slug}:

\begin{lstlisting}[language=json]
{
  "generatedAt": "2026-04-26T05:30:11Z",
  "remotes": [
    { "slug": "login",  "kind": "remote-shell",
      "remoteEntryUrl": "/_remotes/login/remoteEntry.json",
      "exposedModule": "./LoginComponent",
      "displayName": "Login",  "version": "1.0.0", "metadata": {} },
    { "slug": "admin",  "kind": "remote-shell",
      "remoteEntryUrl": "/_remotes/admin/remoteEntry.json",
      "exposedModule": "./AdminShell",
      "displayName": "Admin",  "version": "1.0.0", "metadata": {} },
    /* ...games-three-wizards, games-dodecahedron, wasm-fragment... */
  ]
}
\end{lstlisting}

\section{Shared dependencies}

\snapfile{frontend/federation-base.js} consolidates the shared-package
configuration so the host and all four remotes ship the same Angular
version:

\begin{lstlisting}[language=javascript]
// One source of truth for shareAll + skip helpers.
const STD_SHARED = ['@angular/core','@angular/common','@angular/forms',
                    '@angular/router','rxjs','zone.js'];
const STD_SKIP   = ['@angular/compiler','@angular-architects/native-federation'];

module.exports = {
  snapdragonHost:   () => snapdragonShared({singleton:true,strictVersion:true}),
  snapdragonRemote: () => snapdragonShared({singleton:true,requiredVersion:'auto'}),
};
\end{lstlisting}

The angular-architects builder enforces that \texttt{@angular/core} and
the rest of the standard set are loaded \emph{exactly once} per page;
mismatched versions across remotes produce the
\texttt{NG0203 inject() must be called from an injection context} runtime
error, which is the federation-shaped variant of the more familiar
\texttt{NullInjectorError}.

\section{Routing into a remote}

The \texttt{loadRemote<T>} helper in \snapfile{frontend/src/app/app.routes.ts}
is what every federated route reaches for:

\begin{lstlisting}[language=javascript]
export function loadRemote<T>(
  remoteName: string, exposedModule: string, exportName: string,
): () => Promise<T> {
  return async () => {
    const loadRemoteModule = (globalThis as any)
      .__SNAPDRAGON_LOAD_REMOTE_MODULE__;
    if (!loadRemoteModule) throw new Error('federation runtime missing');
    const m = await loadRemoteModule({ remoteName, exposedModule });
    return m[exportName];
  };
}
\end{lstlisting}

The exported factory is plugged into \texttt{loadComponent} for every
federated route:

\begin{lstlisting}[language=javascript]
{ path: 'login',       loadComponent: loadRemote('login', './LoginComponent', 'LoginComponent') },
{ path: 'admin',       loadComponent: loadRemote('admin', './AdminShell',     'AdminShell')     },
{ path: 'admin/games', loadComponent: loadRemote('admin', './UploadShell',    'UploadShell')    },
\end{lstlisting}

\section{Per-game federation}

A game remote exposes a \emph{Web Component} (built with
\texttt{createCustomElement} and registered with
\texttt{customElements.define}). The host places it as a normal element:

\begin{lstlisting}[language=html]
@switch (activeRemoteSlug()) {
  @case ('three-wizards') {
    <app-three-wizards-remote
      [levelSlug]="activeLevelSlug()"
      [levelId]="activeLevelId()"
      [launchUid]="activeLaunchUid()" />
  }
  @case ('dodecahedron') {
    <app-dodecahedron-remote [launchUid]="activeLaunchUid()" />
  }
}
\end{lstlisting}

The remote internally builds an iframe URL that points at the
\texttt{wasm-fragment} remote:

\begin{lstlisting}[language=javascript]
this.iframeUrl = `/_remotes/wasm-fragment/${this.levelSlug}/index.html` +
                 `?level_id=${this.levelId}&launch_uid=${this.launchUid}` +
                 `&t=${Date.now()}`;
\end{lstlisting}

The cache-busting \texttt{\&t=...} parameter is intentional — the iframe
HTML itself is small and cheap, but the browser is otherwise prone to
caching it across rebuilds and showing a stale Godot bundle.

\section{Error handling}

When a federated chunk fails to load (host nginx down, remote returned
500, network drop), the host's
\snapfile{AppComponent} subscribes to \texttt{Router.events} and shows a
banner with the human-readable label of the broken route, plus a
\textit{Retry} button. The label lookup uses
\texttt{FEDERATED\_ROUTE\_LABELS} (longest-prefix-first, so
\texttt{/admin/games} resolves before \texttt{/admin}).

\section{Maker remotes (RPG + VN)}

Two newer remotes were added alongside the original four to ship the
\emph{maker} side services as federated Angular projects. They follow
exactly the same pattern as \texttt{login} / \texttt{admin}, but their
backing data-source is a separate Spring Boot side service rather than
the BFF:

\begin{itemize}
  \item \texttt{games-rpg-maker} (kind \texttt{remote-shell}) ---
        mounted at \texttt{/rpg-maker}. Talks to
        \snapfile{rpg-builder/} via the BFF's \texttt{/api/rpg/*} proxy
        path.
  \item \texttt{games-vn-maker} (kind \texttt{remote-shell}) ---
        mounted at \texttt{/vn-maker}. Talks to \snapfile{vn-builder/}
        via \texttt{/api/vn/*}.
\end{itemize}

The federation manifest seed for both lives in
\reposrc{backend/src/main/resources/db/changelog/037-rpg-and-vn-maker-federation.xml}
so the host shell discovers them at bootstrap without a redeploy. The
host nginx config (and the TLS twin --- see chapter~\ref{ch:builddeploy})
needs matching \texttt{location /\_remotes/games-rpg-maker/}
and \texttt{/\_remotes/games-vn-maker/} proxy blocks alongside the
\texttt{/api/rpg/} and \texttt{/api/vn/} sub-service proxies; drift
between the plain and TLS variants produces ``Couldn't reach \emph{X}
--- federated module didn't load.''

\begin{snapcaution}[title={CORP saves WASM iframes}]
WASM games loading inside the federated iframe need
\texttt{Cross-Origin-Resource-Policy: cross-origin} on the
\texttt{wasm-fragment} response \emph{and} the host nginx must strip any
duplicate COOP/COEP headers from the upstream and re-emit a single
canonical set. Missing CORP shows up as
\texttt{net::ERR\_BLOCKED\_BY\_RESPONSE}; duplicates show up as
\texttt{NotSameOriginAfterDefaultedToSameOriginByCoep}. Both are silent
in DevTools' \emph{Console} but show clearly in \emph{Network}'s
\textit{Cross-origin isolation} column.
\end{snapcaution}

\chapter{Data model}\label{ch:datamodel}

PostgreSQL holds the canonical relational state. Cassandra is a
write-through cross-pod cache for movement events, an audit log for game
state transitions, and a session store for short-lived browser tokens.
Liquibase drives schema migration through 24 changelogs at the head of
the federation branch; only Postgres is migrated this way — Cassandra
schemas are bootstrapped by an init container.

\section{Core relational schema}

\begin{figure}[!h]
\centering
\snapdiagram{diagrams/data-model.tex}
\caption{Core schema clusters. Arrows are foreign keys; tables in the
\enquote{federation \& access} cluster were added by migrations 018–024 and
are the post-federation additions to a previously single-game schema.}
\end{figure}

\subsection{Games and levels}

\begin{tabularx}{\linewidth}{l X}
\toprule
\textbf{Table} & \textbf{Purpose} \\
\midrule
\texttt{games} & Umbrella row per game.
\texttt{slug} is the federation key (matches
\texttt{federation\_modules.slug}); \texttt{current\_version\_dir} is the
hash-prefixed directory under \snapfile{app.uploads.games-dir} that the
runtime serves from. \\
\texttt{levels} & One row per WASM bundle. \texttt{asset\_url} resolves to
either \texttt{/api/assets/levels/<slug>/index.html} (legacy) or
\texttt{/\_remotes/wasm-fragment/<slug>/index.html} (federated). \\
\texttt{scene\_elements} & Per-level DTO rows that
\snapfile{scene_loader.gd} pulls from
\texttt{GET /api/levels/\{id\}/scene}. \\
\texttt{level\_assets} & Optional registry for runtime-loaded textures /
plugins / audio, looked up by \texttt{label} from GDScript via
\texttt{ApiClient.resolve\_asset(level\_id, label, cb)}. \\
\texttt{level\_prerequisites} & Quest gating: a level shows as locked
on the home page until its quest is complete. \\
\bottomrule
\end{tabularx}

\subsection{Federation, access, and uploads}

\begin{tabularx}{\linewidth}{l X}
\toprule
\textbf{Table} & \textbf{Purpose} \\
\midrule
\texttt{federation\_modules} & Source of truth for the federation
manifest. Rows include \texttt{kind}
(\texttt{remote-shell|remote-game|fragment}),
\texttt{remote\_entry\_url}, and \texttt{exposed\_module}. Read by
\snapfile{FederationController}. \\
\texttt{game\_access} & Per-game visibility policy.
\texttt{(scope, scope\_value)} can be
\texttt{(all, NULL)} / \texttt{(role, ROLE\_PLAY\_X)} / \texttt{(user, <casdoor-id>)};
the unique constraint \texttt{uq\_game\_access\_triple} keeps inserts
idempotent. \\
\texttt{game\_launches} & 256-bit URL-safe \texttt{launch\_uid} per
session, with a \texttt{version\_dir} snapshot from
\texttt{games.current\_version\_dir} at mint time so an in-flight launch
keeps reading from its install even after a hot upload supersedes it.
TTL 15 minutes; the \snapfile{GameLaunchSweeper} closes stale rows. \\
\texttt{game\_uploads} & One row per upload zip, with the
\texttt{manifest\_json} text (jsonb $\rightarrow$ text per migration
024-1) and a \texttt{superseded\_at}
timestamp on rollback. Bound by
\texttt{app.uploads.versions-to-keep} via
\snapfile{GameUploadPruner}. \\
\bottomrule
\end{tabularx}

\subsection{Players, sessions, and progression}

\begin{tabularx}{\linewidth}{l X}
\toprule
\textbf{Table} & \textbf{Purpose} \\
\midrule
\texttt{characters} & One canonical character per Casdoor user, with
\texttt{position\_x/y/z}, \texttt{rotation\_y}, HP/MP/XP, and
\texttt{current\_level\_id}. \\
\texttt{spell\_loadout} & Per-character spell-slot binding (key index
$\rightarrow$ spell). \\
\texttt{quests}, \texttt{character\_quests} & Quest catalog and
per-character progress. \\
\texttt{items}, \texttt{character\_items} & Inventory. \\
\texttt{controls} & Sessionised input state — at most one row per
\texttt{(character\_id, session\_id)}; only the holder of the controller
role for a session may POST. \\
\texttt{player\_positions} & One row per
\texttt{(character\_id, level\_id)} — Liquibase 027 promoted the
legacy single-row-per-character layout into a per-level layout so a
character traversing cave $\rightarrow$ field $\rightarrow$ cave
returns to the cave coordinates they last had on re-entry instead of
clobbering them across levels. Cassandra's
\texttt{character\_position\_cache} (composite partition
\texttt{(character\_id, level\_id)}) is the cross-pod read cache that
lets pod B see a position pod A wrote within one SSE round-trip;
Postgres backs the cache when the 5-minute TTL drops a row. \\
\bottomrule
\end{tabularx}

\section{Migration ordering}

Migrations 001–017 are pre-federation. The federation work landed in:

\begin{itemize}
  \item \reposrc{backend/src/main/resources/db/changelog/018-federation-modules.xml} —
        \texttt{federation\_modules} + seed rows for the four remotes plus
        \texttt{wasm-fragment}.
  \item \reposrc{backend/src/main/resources/db/changelog/019-game-access.xml} —
        \texttt{game\_access} + \texttt{uq\_game\_access\_triple}.
  \item \reposrc{backend/src/main/resources/db/changelog/020-game-launches.xml} —
        \texttt{game\_launches}.
  \item \reposrc{backend/src/main/resources/db/changelog/021-game-uploads.xml} —
        \texttt{game\_uploads}.
  \item \reposrc{backend/src/main/resources/db/changelog/022-federation-relative-urls.xml} —
        rewrites absolute remote URLs to relative paths.
  \item \reposrc{backend/src/main/resources/db/changelog/023-launch-version-pin.xml} —
        adds \texttt{version\_dir} to \texttt{game\_launches}.
  \item \reposrc{backend/src/main/resources/db/changelog/024-game-uploads-manifest-text.xml} —
        \texttt{ALTER manifest\_json} jsonb $\rightarrow$ text plus
        \texttt{setval(games\_id\_seq, MAX(id))} (the integration test
        surfaced both as latent bugs).
\end{itemize}

\section{R2DBC + JDBC coexistence}

Runtime data access is non-blocking R2DBC. Liquibase needs a blocking
JDBC \texttt{DataSource}, so \snapfile{JdbcDataSourceConfig} declares an
explicit \texttt{@Primary HikariDataSource} (pool size 3) just for the
migration phase. R2DBC then connects in parallel against the same
database via the \texttt{spring.r2dbc.url} property.

\begin{snapcaution}[title={MVCC snapshot race}]
\texttt{@PostConstruct} services that touch the database can run
\emph{before} Liquibase commits, locking their R2DBC connection into a
pre-migration snapshot. The fix in this codebase is \texttt{@DependsOn(\"liquibase\")}
on \snapfile{GameTimeService} and \snapfile{GameStateService}, plus
\texttt{spring.r2dbc.pool.initial-size=0} in tests so the pool only opens
connections after Liquibase finishes. Without these, the symptom is
\texttt{BadSqlGrammarException: relation does not exist} from a query
the test fixture clearly sets up.
\end{snapcaution}

\section{Cassandra cache}

The Cassandra keyspace holds:

\begin{itemize}
  \item \textbf{Movement events} — write-through buffer for
        \texttt{ControlEvent}s; the SSE stream reads from Postgres but
        Cassandra holds the recent N events for fast cross-pod replay.
  \item \textbf{Audit log} — every state transition emitted by
        \snapfile{GameStateService} is persisted with
        \texttt{(character\_id, occurred\_at)} as the primary key.
  \item \textbf{Session store} — short-lived per-tab tokens bound to a
        Casdoor user, used by the BFF to refuse \texttt{/api/controls}
        requests from a tab that's not currently the controller for the
        target session.
\end{itemize}

A nightly \snapfile{CassandraSyncService} reconciles audit-log totals
back into Postgres counters so dashboards stay consistent across reboots.

\section{Generalised input event surface (Liquibase 026)}

The five legacy game tables (\texttt{characters}, \texttt{controls},
\texttt{spell\_loadout}, ...) cover JRPG-shaped gameplay; the input
surface added in changeset 026 covers \emph{everything Godot can
consume}: gamepad button / axis, mouse motion / button / wheel,
keyboard up / down, touch gestures, audio cues, world events,
plus the legacy control commands as one kind among many.

Three reference tables anchor the schema; an append-only audit
table writes against them.

\begin{tabularx}{\linewidth}{l X}
\toprule
\textbf{Table} & \textbf{Purpose} \\
\midrule
\texttt{input\_event\_kinds} & Reference. Primary key on
\texttt{kind}. Closed vocabulary of input categories: \texttt{gamepad},
\texttt{mouse}, \texttt{keyboard}, \texttt{touch}, \texttt{audio},
\texttt{world}, \texttt{control}. Seeded by Liquibase 026. \\

\texttt{input\_event\_types} & Reference. \textbf{Composite}
primary key \texttt{(kind, type)}. Per-kind type vocabulary
(\texttt{gamepad: button\_down / button\_up / axis / connected /
disconnected / rumble\_ack}, \texttt{mouse: motion / button\_down /
button\_up / wheel}, ...) plus an \texttt{example\_payload} (jsonb)
used by the upload-time validator and the OpenAPI explorer's sample
generator. \\

\texttt{game\_event\_support} & Per-game opt-in. \textbf{Composite}
primary key \texttt{(game\_slug, event\_kind, event\_type)}, FK to
\texttt{input\_event\_types(kind, type)}. The Snapdragon Export
plugin writes one row per declared event during upload; the BFF
refuses any input emit whose triple isn't enabled here, regardless
of what the client sends. \\

\texttt{input\_events} & Append-only audit log. FK on
\texttt{(event\_kind, event\_type)} into
\texttt{input\_event\_types} so a row can never reference an
unknown type. Indexed on \texttt{(character\_id, occurred\_at)},
\texttt{(session\_id, occurred\_at)}, and
\texttt{(event\_kind, occurred\_at)} for replay + analytics. \\

\texttt{games.features} & New \texttt{jsonb} column. Per-game
feature flags: \texttt{spell\_loadout}, \texttt{quests},
\texttt{items}, \texttt{char\_progression} (the JRPG subsystems —
non-JRPG games turn these off and the matching endpoints return
\texttt{405}); \texttt{audio}, \texttt{gamepad}, \texttt{mouse},
\texttt{keyboard}, \texttt{touch} (input-kind gates — a game with
\texttt{gamepad: false} can't accidentally write gamepad rows). \\
\bottomrule
\end{tabularx}

\bigskip

The validation flow on every input emit is three gates:
\begin{enumerate}
  \item \texttt{(kind, type)} must exist in
        \texttt{input\_event\_types} (400 if not).
  \item \texttt{games.features.<kind>} must be true
        for the active game (405 if not).
  \item \texttt{(game\_slug, kind, type)} must be present in
        \texttt{game\_event\_support} (405 if not).
\end{enumerate}

Same surface persists every browser-tier gamepad poll, every
WASM-iframe console line forwarded by the host shell, every
audio cue the Godot bundle emits. The Rust GDExtension's
\texttt{HostBridge} ships matching typed signals
(\texttt{gamepad\_received}, \texttt{mouse\_received},
\texttt{keyboard\_received}, \texttt{touch\_received},
\texttt{audio\_received}, plus outbound \texttt{audio\_egress} and
\texttt{console\_egress}) so a Godot developer wires gameplay to
named signals without caring whether the source was a player's
USB Pro Controller, a \texttt{<kbd>W</kbd>} keypress, or an SSE
emit from a cross-tab subscriber.

\chapter{API reference}\label{ch:api}

This chapter is the printed companion to the generated Javadoc HTML at
\reposrc{docs/publishing/javadoc-output/index.html}. Every controller is
listed once with its endpoints, auth posture, and a one-line description
of the contract; the Javadoc has the per-method signatures, parameter
descriptions, and full HTTP examples. The Javadoc is rebuilt by
\snapcommand{make docs-javadocs} and is themed to match this document
(see \snapfile{docs/publishing/javadoc-theme/snapdragon.css}).

\bigskip

The reactive controllers below carry a \emph{marble diagram} showing
their Project Reactor chain. Marbles use distinct shapes for distinct
runtime types so the topology is readable at a glance:
\snapmarbleshape{marbleMono}~\texttt{Mono<T>} value,
\snapmarbleshape{marbleFlux}~\texttt{Flux<T>} element,
\snapmarbleshape{marbleSSE}~\texttt{ServerSentEvent} (HTTP-tier emission),
\snapmarbleshape{marbleTx}~persistence-tier value (R2DBC row / post-tx),
\snapmarbleshape{marbleSig}~control / signal,
\snapmarbleshape{marbleErr}~error,
\snapmarbleshape{marbleDone}~\texttt{onComplete}. Each diagram's legend
lists the carried object's relevant fields and flags so a reader can
cross-reference the chain with the BFF source code.

\section{Public surface}

\subsection{\texttt{AuthController} — \texttt{/api/auth}}
Public reverse-proxy of the Casdoor login flow.

\begin{itemize}
  \item \texttt{POST /api/auth/token} — exchanges a Casdoor authorization
        code for an access token; the federated \texttt{login} remote
        calls this directly.
  \item \texttt{GET /api/auth/me} — returns the resolved Casdoor user
        from the bearer token.
\end{itemize}

\snapdiagram{marbles/auth-controller.tex}
\snapmarblelegend{
  \snapmarbleshape{marbleMono} & \texttt{TokenExchangeRequest}
    & \texttt{code} (Casdoor authorization code) {\rm $\cdot$} \texttt{redirectUri} {\rm $\cdot$} \texttt{state} (CSRF echo). \\
  \snapmarbleshape{marbleSSE}  & Upstream JSON
    & raw Casdoor response: \texttt{access\_token}, \texttt{token\_type}, \texttt{expires\_in}, \texttt{refresh\_token}. \\
  \snapmarbleshape{marbleMono} & \texttt{TokenResponse}
    & decoded DTO: \texttt{accessToken}, \texttt{expiresAt} (Instant), \texttt{refreshToken}, \texttt{casdoorUserId}. \\
  \snapmarbleshape{marbleDone} & \texttt{onComplete}
    & terminal: a single value plus completion. The host writes it to \texttt{sessionStorage.tw\_auth}. \\
}

\subsection{\texttt{GameStreamController} — \texttt{/api/game}}
Public SSE merge stream.

\begin{itemize}
  \item \texttt{GET /api/game/stream} — returns five typed event channels
        (\texttt{movement}, \texttt{world}, \texttt{character},
        \texttt{time}, \texttt{scene}) plus a 15-second keepalive.
\end{itemize}

\snapdiagram{marbles/game-stream-controller.tex}
\snapmarblelegend{
  \snapmarbleshape{marbleFlux} & \texttt{PlayerMovement}
    & \texttt{characterId} {\rm $\cdot$} \texttt{x,y,z} {\rm $\cdot$} \texttt{rotationY} {\rm $\cdot$} \texttt{action} {\rm $\cdot$} \texttt{pressed} (bool) {\rm $\cdot$} \texttt{tick}. \\
  \snapmarbleshape{marbleFlux} & \texttt{WorldEvent}
    & \texttt{type} (\texttt{spawn}/\texttt{despawn}/\texttt{transition}) {\rm $\cdot$} \texttt{payload} (Map). \\
  \snapmarbleshape{marbleFlux} & \texttt{CharacterEvent}
    & \texttt{characterId} {\rm $\cdot$} \texttt{hp}, \texttt{mp}, \texttt{xp} {\rm $\cdot$} \texttt{flags}: \texttt{stunned}, \texttt{burning}, \texttt{frozen}. \\
  \snapmarbleshape{marbleFlux} & \texttt{GameTimeTick}
    & \texttt{day} {\rm $\cdot$} \texttt{hour} {\rm $\cdot$} \texttt{minute} {\rm $\cdot$} \texttt{phase} (\texttt{day}/\texttt{dusk}/\texttt{night}/\texttt{dawn}). \\
  \snapmarbleshape{marbleSig}  & ping (15-second keepalive)
    & no payload — fires the moment the subscribe lands so \texttt{onopen} flushes. \\
  \snapmarbleshape{marbleSSE}  & \texttt{ServerSentEvent}
    & merged downstream: \texttt{event} (the channel name) {\rm $\cdot$} \texttt{data} (JSON of the source value) {\rm $\cdot$} \texttt{id} (sequence). \\
}

\subsection{\texttt{GameController} — \texttt{/api/games}}
Public game catalog (filtered by \texttt{game\_access} when a JWT is
present).

\begin{itemize}
  \item \texttt{GET /api/games} — list of visible games.
  \item \texttt{GET /api/games/\{slug\}} — one game by slug.
\end{itemize}

\subsection{\texttt{LevelController} — \texttt{/api/levels}}
Public per-level metadata and scene DTOs.

\begin{itemize}
  \item \texttt{GET /api/levels/\{id\}} — level row including
        \texttt{asset\_url}.
  \item \texttt{GET /api/levels/\{id\}/scene} — typed
        \texttt{SceneDto} that \snapfile{scene_loader.gd} consumes.
\end{itemize}

\snapdiagram{marbles/level-controller.tex}
\snapmarblelegend{
  \snapmarbleshape{marbleMono} & \texttt{Level} (catalog row)
    & \texttt{id} {\rm $\cdot$} \texttt{slug} {\rm $\cdot$} \texttt{sceneType} {\rm $\cdot$} \texttt{assetUrl} {\rm $\cdot$} \texttt{gameId} (FK) {\rm $\cdot$} \texttt{sceneConfig} (JSON). \\
  \snapmarbleshape{marbleTx}   & \texttt{SceneElement} (one DB row per spawn)
    & \texttt{elementType} (\texttt{fairy}/\texttt{torch}/...) {\rm $\cdot$} \texttt{positionX,Y,Z} {\rm $\cdot$} \texttt{properties} (JSON) {\rm $\cdot$} \texttt{sortOrder}. \\
  \snapmarbleshape{marbleMono} & \texttt{SceneDto}
    & \texttt{levelId} {\rm $\cdot$} \texttt{sceneType} {\rm $\cdot$} \texttt{elements: List<ElementDto>} (the Godot \texttt{scene\_loader} reads this directly). \\
}

\subsection{\texttt{LevelAssetController} — \texttt{/api/levels/\{id\}/assets}}
Public per-level asset registry — \texttt{POST} to register, \texttt{GET}
to list, \texttt{GET /\{label\}} to resolve.

\subsection{\texttt{AssetController} — \texttt{/api/assets}}
Public legacy classpath asset serving (\texttt{/api/assets/levels/<slug>/*}).
Honours \texttt{app.assets.require-launch-uid=true} for defense in depth.

\subsection{\texttt{ControlController} — \texttt{/api/controls}}
Session-validated control event ingestion. The session check is enforced
through the JWT subject; the legacy \texttt{X-Casdoor-Id} header path is
gone.

\snapdiagram{marbles/control-controller.tex}
\snapmarblelegend{
  \snapmarbleshape{marbleSig}  & \texttt{ControlEvent}
    & \texttt{characterId} (path var) {\rm $\cdot$} \texttt{sessionId} (header) {\rm $\cdot$} \texttt{action} (\texttt{move}/\texttt{cast}/\texttt{interact}) {\rm $\cdot$} \texttt{x,y} {\rm $\cdot$} \texttt{pressed} (bool) {\rm $\cdot$} \texttt{spellSlot}? . \\
  \snapmarbleshape{marbleTx}   & persisted row
    & \texttt{controls} table: same fields plus \texttt{persistedAt} (Instant) and the unique \texttt{(characterId, sessionId)} constraint. \\
  \snapmarbleshape{marbleFlux} & sink emission
    & a copy of the persisted row pushed onto \texttt{GameStateService.movementSink} so every connected SSE subscriber sees it. \\
  \snapmarbleshape{marbleErr}  & \texttt{Forbidden403}
    & raised by \texttt{filterWhen} when the JWT subject isn't the controller-role holder for that session; never reaches the sink. \\
}

\subsection{\texttt{FederationController} — \texttt{/api/federation}}
Public federation manifest.

\begin{itemize}
  \item \texttt{GET /api/federation/manifest} — DB-backed manifest read
        from \texttt{federation\_modules}, ordered by
        \texttt{kind}/\texttt{slug}.
\end{itemize}

\snapdiagram{marbles/federation-controller.tex}
\snapmarblelegend{
  \snapmarbleshape{marbleTx}   & \texttt{FederationModule} (Postgres row)
    & \texttt{slug} {\rm $\cdot$} \texttt{kind} (\texttt{remote-shell} / \texttt{remote-game} / \texttt{fragment}) {\rm $\cdot$} \texttt{remoteEntryUrl} {\rm $\cdot$} \texttt{exposedModule} {\rm $\cdot$} \texttt{enabled} (bool). \\
  \snapmarbleshape{marbleMono} & \texttt{FederationRemoteDto}
    & + \texttt{displayName} {\rm $\cdot$} \texttt{version} {\rm $\cdot$} \texttt{metadata} (Map); ordered by \texttt{kind, slug}. \\
  \snapmarbleshape{marbleMono} & \texttt{ManifestResponse}
    & \texttt{generatedAt} (Instant, ISO-8601) {\rm $\cdot$} \texttt{remotes: List<FederationRemoteDto>}; consumed by the host shell BEFORE \texttt{bootstrapApplication}. \\
}

\section{Authenticated surface}

\subsection{\texttt{SessionController} — \texttt{/api/sessions}}
Session connect/disconnect; \texttt{POST /api/sessions/connect} claims
the controller role if no one else holds it.

\subsection{\texttt{CharacterController} — \texttt{/api/characters}}
Character read/update including \texttt{POST /api/characters/\{id\}/move}
which emits a \texttt{movement} SSE event after persisting.

\subsection{\texttt{SkillController} — \texttt{/api/skills}}
Skill catalog and per-character unlock state.

\subsection{\texttt{ItemController} — \texttt{/api/items}}
Inventory + per-character item state.

\subsection{\texttt{QuestController} — \texttt{/api/quests}}
Quest catalog and per-character progress.

\subsection{\texttt{GameTimeController} — \texttt{/api/time}}
Authenticated game-time read; the public SSE stream emits the same data
on the \texttt{time} channel.

\subsection{\texttt{GameLaunchController} — \texttt{/api/games}, \texttt{/api/admin/launches}}

\begin{itemize}
  \item \texttt{POST /api/games/\{slug\}/launch} — mints a
        \texttt{launch\_uid}, snapshots
        \texttt{games.current\_version\_dir}, returns
        \texttt{\{uid, expiresAt, gameSlug\}}.
  \item \texttt{DELETE /api/games/launches/\{uid\}} — owner
        self-terminate.
  \item \texttt{DELETE /api/admin/launches/\{uid\}} — admin
        force-terminate.
\end{itemize}

\snapdiagram{marbles/game-launch-controller.tex}
\snapmarblelegend{
  \snapmarbleshape{marbleMono} & \texttt{LaunchPayload}
    & \texttt{slug} (path var) {\rm $\cdot$} \texttt{casdoorUserId} (JWT) {\rm $\cdot$} \texttt{sessionId} (header). \\
  \snapmarbleshape{marbleMono} & \texttt{LaunchDraft}
    & + \texttt{uid} (256-bit URL-safe) {\rm $\cdot$} \texttt{expiresAt = now + 15 min} {\rm $\cdot$} \texttt{version\_dir = NULL} (yet). \\
  \snapmarbleshape{marbleMono} & version-pinned draft
    & \texttt{LaunchDraft} merged with \texttt{games.current\_version\_dir} so the row is reproducible after a hot upload. \\
  \snapmarbleshape{marbleTx}   & \texttt{LaunchResponse}
    & post-commit row: \texttt{uid} {\rm $\cdot$} \texttt{gameSlug} {\rm $\cdot$} \texttt{expiresAt} {\rm $\cdot$} \texttt{versionDir}; threaded into the iframe URL as \texttt{?launch\_uid=...} . \\
}

\section{Admin-only surface}

\subsection{\texttt{AdminController} — \texttt{/api/admin}}
Admin dashboard data: per-user stats, active sessions, recent uploads.

\subsection{\texttt{GameUploadController} — \texttt{/api/admin/games}}

\begin{itemize}
  \item \texttt{POST /api/admin/games/upload} — multipart bundle zip;
        the install pipeline validates manifest schema v2 + slug regex +
        zip-slip-safe paths, hashes the bundle, atomically moves into
        \texttt{<gamesDir>/<slug>/<sha-prefix>}, runs the bundle's
        \texttt{register.sql} inside a single R2DBC transaction, and
        supersedes prior uploads. Returns
        \texttt{\{slug, contentHash, versionDir, uploadedAt\}}.
  \item \texttt{GET /api/admin/games/\{slug\}/uploads} — rollback history,
        newest first; \texttt{supersededAt = NULL} marks the current.
\end{itemize}

\snapdiagram{marbles/game-upload-controller.tex}
\snapmarblelegend{
  \snapmarbleshape{marbleMono} & \texttt{UploadRequest}
    & \texttt{file} (MultipartFile, the zip) {\rm $\cdot$} \texttt{casdoorUserId} (from JWT) {\rm $\cdot$} \texttt{contentLength}. \\
  \snapmarbleshape{marbleMono} & \texttt{ValidatedZip}
    & \texttt{manifestJson} (schema\_v=2) {\rm $\cdot$} \texttt{slug} (regex \texttt{[a-z0-9-]+}) {\rm $\cdot$} \texttt{registerSql} {\rm $\cdot$} \texttt{flag}: \texttt{zipSlipSafe=true}. \\
  \snapmarbleshape{marbleMono} & \texttt{StagedBundle}
    & \texttt{slug} {\rm $\cdot$} \texttt{contentHash} (SHA-256, 64-char hex) {\rm $\cdot$} \texttt{stagedDir} (\texttt{<gamesDir>/<slug>/<sha-prefix>}). \\
  \snapmarbleshape{marbleTx}   & \texttt{InstalledBundle} (post-commit)
    & \texttt{slug} {\rm $\cdot$} \texttt{contentHash} {\rm $\cdot$} \texttt{versionDir} {\rm $\cdot$} \texttt{uploadedAt} {\rm $\cdot$} \texttt{supersededAt = NULL}. \\
  \snapmarbleshape{marbleErr}  & \texttt{UploadException}
    & \texttt{status} (\texttt{400}/\texttt{409}/\texttt{500}) {\rm $\cdot$} \texttt{reason} {\rm $\cdot$} \texttt{slug} (if known); rolls back the tx + on-disk move. \\
  \snapmarbleshape{marbleDone} & \texttt{onComplete}
    & terminal: response carries the four \texttt{InstalledBundle} fields and HTTP 200. \\
}

\subsection{\texttt{GameAccessController} — \texttt{/api/admin/games}}

\begin{itemize}
  \item \texttt{POST /api/admin/games/\{slug\}/access} — body
        \texttt{\{scope, value?\}}; absorbs duplicates via
        \texttt{uq\_game\_access\_triple}.
  \item \texttt{DELETE /api/admin/games/\{slug\}/access?scope=\&value=}
        — removes a row.
  \item \texttt{GET /api/admin/games/\{slug\}/access} — additive list of
        active rows.
\end{itemize}

\subsection{\texttt{RoleProxyController} — \texttt{/api/admin/roles}}
Server-side proxy of \texttt{/api/get-roles?owner=threewizards} so the
admin UI's \enquote{grant role} dropdown loads from the live Casdoor
catalog without exposing the Casdoor admin client secret to the browser.

\section{Where to find more detail}

The full Javadoc is published to
\reposrc{docs/publishing/javadoc-output/} when \snapcommand{make docs-javadocs}
runs. The CSS at \reposrc{docs/publishing/javadoc-theme/snapdragon.css}
themes the output to match this document — same accent purple, same
\texttt{Inter}/\texttt{Segoe UI} type stack, same code-block treatment.
The Javadoc HTML is git-ignored.

\chapter{Build and deploy}

The repository's root \snapfile{Makefile} is the canonical entry point.
Every target is idempotent and safe to re-run; the build script for this
document (\reposrc{docs/publishing/Makefile}) follows the same convention.

\section{Build chain}

\begin{figure}[!h]
\centering
\snapdiagram{diagrams/build-chain.tex}
\caption{Build chain. Each rectangle is a Make target; arrows are
dependencies. The \texttt{docs-*} targets at the right are the additions
introduced by this document's build pipeline.}
\end{figure}

\subsection{Native (Linux) build}

\begin{lstlisting}[language=bash]
make build-rust       # cargo build --release + (optional) wasm32-unknown-emscripten
make build-godot      # exports every preset in every game project under games/<dir>/
make build-backend    # ./mvnw package -DskipTests
make build-frontend   # ng build for host + four remotes
make build            # runs all of the above in dependency order
\end{lstlisting}

The \texttt{build-rust} target builds the native \snapfile{.so} for the
editor and tries the WASM target; if \texttt{emcc} isn't on \texttt{PATH}
the WASM step prints a warning and continues — the rest of the build
chain doesn't depend on it for normal browser play through Godot's WASM
target.

\subsection{Docker compose (recommended)}

\begin{lstlisting}[language=bash]
docker compose up -d                  # postgres + cassandra + casdoor + bff + frontend + remotes
make e2e                              # full Playwright run against the live stack
make e2e-native                       # same, against the GraalVM native image overlay
\end{lstlisting}

\begin{snapcaution}[title={Docker 29.x needs buildx}]
Docker 29.x removed the legacy in-daemon image builder. \texttt{docker
compose build} will hang silently if \texttt{docker-buildx} isn't
installed. The fix is a user-local install of the v0.18 binary into
\texttt{\textasciitilde/.docker/cli-plugins/docker-buildx} (no sudo
needed). The README's \emph{Running with Docker Compose} section
documents the exact command.
\end{snapcaution}

\section{GraalVM native image}

The BFF compiles to a native binary via GraalVM \texttt{native-image}.
The driver is a separate build profile in \snapfile{backend/pom.xml},
plus a Dockerfile (\snapfile{backend/Dockerfile.native}) and a compose
overlay (\snapfile{compose.native.yml}).

\begin{lstlisting}[language=bash]
make build-backend-native     # mvn -Pnative native:compile (~6 minutes on a workstation)
make e2e-native               # docker compose -f compose.yml -f compose.native.yml up -d
\end{lstlisting}

The native build needs reflection hints registered for every R2DBC
mapping class, every controller DTO, and every Liquibase changelog
entity. \snapfile{R2dbcConvertersConfig} centralises these so adding a
new entity touches one file.

\section{The publishing pipeline}

This document, the Javadoc, and the three slide decks are produced by
\snapfile{docs/publishing/}. The published artefacts are git-ignored.

\begin{lstlisting}[language=bash]
make docs-pdf          # docs/publishing/architecture/build/architecture.pdf
make docs-javadocs     # docs/publishing/javadoc-output/index.html
make docs-slides       # docs/publishing/slides/build/*.pptx
make docs-all          # all of the above
make docs-clean        # remove every produced artefact
\end{lstlisting}

\subsection{LaTeX prerequisites}

\begin{itemize}
  \item TeX Live with \texttt{pdflatex}, \texttt{xelatex},
        \texttt{tcolorbox}, \texttt{tikz}, \texttt{titlesec},
        \texttt{titletoc}, \texttt{enumitem}, \texttt{fontawesome},
        \texttt{inconsolata}.
  \item On Debian / Ubuntu: \texttt{apt install texlive-latex-extra
        texlive-fonts-extra texlive-pictures}.
  \item On Gentoo: the \texttt{tex-live} portage package with USE
        \texttt{extra} produces the same set.
\end{itemize}

\subsection{Javadoc theming}

\snapfile{backend/pom.xml} configures the Maven Javadoc plugin to read
\snapfile{docs/publishing/javadoc-theme/snapdragon.css} as the override
stylesheet and \snapfile{docs/publishing/javadoc-theme/header.html} as
the page header. The override is purely additive — the standard doclet's
default stylesheet still loads first; the Snapdragon CSS layers on top.

\subsection{PowerPoint generator}

The slide decks are generated by Python scripts under
\reposrc{docs/publishing/slides/}. Each script imports
\texttt{python-pptx}, builds the deck programmatically, and writes the
\texttt{.pptx} file to \snapfile{docs/publishing/slides/build/}. The
scripts are deliberately written so a non-Python audience can edit slide
content by editing the data structure at the top of each file —
template placeholders are not used.

\section{CI hooks}

A pre-push hook is shipped at
\reposrc{docs/publishing/scripts/install-hooks.sh}. It installs
\snapfile{.git/hooks/pre-push} that re-runs \texttt{make docs-pdf} on
any push that touches \snapfile{docs/publishing/} or
\snapfile{docs/PLAN-module-federation-and-native.md}. Failures abort the
push so a broken build can't land on \texttt{main}. The hook is opt-in;
the script is idempotent and prints what it changed.

\chapter{Extending the portal}

This chapter is the load-bearing summary of the three slide decks under
\reposrc{docs/publishing/slides/}: how to add a new game, a new role, a
new federated route, or a new admin view. The slides have screenshots
and click-by-click walkthroughs; this chapter is the canonical
\emph{contract} those walkthroughs follow.

\section{Adding a new game}

There are two paths.

\subsection{Path A — upload at runtime}

Every supported game ships an editor plugin
(\snapfile{addons/snapdragon_export/}) under
\reposrc{games/three-wizards-in-a-trenchcoat/} that produces a deploy-ready
zip with a v2 \texttt{manifest.json}, a \texttt{register.sql} insert
script, and the \texttt{bundle/} contents:

\begin{enumerate}
  \item In the Godot editor, \emph{Project~$\rightarrow$~Tools~$\rightarrow$~Snapdragon~$\rightarrow$~Export\dots}.
  \item Fill in slug, display name, description, scene type. The plugin
        defaults the federation \texttt{remote\_entry\_url} to
        \texttt{/\_remotes/games-<slug>/remoteEntry.json}; if you don't
        plan to ship a per-game frontend remote, change it to
        \texttt{/\_remotes/wasm-fragment/remoteEntry.json} so the host
        loads the bundle through the shared fragment.
  \item Click \emph{Export}. The plugin writes a zip to your chosen
        path.
  \item Hit \texttt{POST /api/admin/games/upload} with the zip — the
        admin UI's \emph{Upload game bundle} card does this for you.
        The install transaction validates, hashes, atomically moves the
        bundle, and runs \texttt{register.sql} which inserts into
        \texttt{games}, \texttt{federation\_modules}, and
        \texttt{game\_access} all in one Postgres transaction.
\end{enumerate}

The host shell discovers the new game on the next federation manifest
fetch (currently every 60s; see \snapfile{FederationService}). No host
redeploy is required.

\subsection{Path B — commit the game}

For a game you want to ship in the repository:

\begin{enumerate}
  \item Drop the Godot project under
        \snapfile{games/<game-dir>/project.godot}, with its own
        \snapfile{export_presets.cfg} and
        \snapfile{addons/snapdragon_export/}.
  \item Add the directory to \texttt{GAME\_PROJECTS} and define
        \texttt{PRESETS\_<dir>} in the root \snapfile{Makefile}.
  \item Add a Liquibase changeset that inserts the \texttt{games} +
        \texttt{federation\_modules} + \texttt{game\_access} rows. See
        \reposrc{backend/src/main/resources/db/changelog/017-snapdragon-catalog.xml}
        for the canonical pattern.
  \item Optionally add a \texttt{games-<slug>} service to
        \snapfile{docker-compose.yml} if the game ships its own
        Module Federation remote (separate Angular project under
        \snapfile{frontend/projects/games-<slug>/}).
  \item Run \texttt{make build} and \texttt{docker compose up -d}.
\end{enumerate}

\section{Adding a new federated route}

\begin{enumerate}
  \item Build the route's component as a standalone Angular project
        under \snapfile{frontend/projects/<remote-name>/}, with its own
        \snapfile{federation.config.js} extending
        \snapfile{frontend/federation-base.js}.
  \item Expose the entry component in
        \snapfile{frontend/projects/<remote-name>/federation.config.js}'s
        \texttt{exposes} map.
  \item Add a row to \texttt{federation\_modules} (Liquibase changeset
        recommended). Set \texttt{kind} to \texttt{remote-shell} for a
        full route component.
  \item In \snapfile{frontend/src/app/app.routes.ts}, plug
        \texttt{loadRemote('<remote-name>', './<ExposedSymbol>',
        '<ExposedSymbol>')} into a \texttt{loadComponent} entry.
  \item Add an nginx \texttt{location /\_remotes/<remote-name>/} block
        in \snapfile{frontend/nginx.conf} pointing at the per-remote
        compose service.
  \item Add the per-remote service to \snapfile{docker-compose.yml}.
\end{enumerate}

\section{Adding a new admin view}

The admin remote (\reposrc{frontend/projects/admin/}) exposes any
top-level standalone component listed in its \texttt{exposes} map. To
add a view:

\begin{enumerate}
  \item Add a standalone component under
        \snapfile{frontend/projects/admin/src/app/}.
  \item Add it to the \texttt{exposes} map in
        \snapfile{frontend/projects/admin/federation.config.js}.
  \item Wire a route in the host's
        \snapfile{frontend/src/app/app.routes.ts} pointing at the
        exposed name.
\end{enumerate}

DTOs cannot be shared across the federation boundary at runtime
(TypeScript types erase). Redeclare interfaces inside the remote — see
the existing pattern in
\reposrc{frontend/projects/admin/src/app/upload-remote.component.ts}.

\section{Adding a new role}

Casdoor is the source of truth for the role catalog. The portal
enforces visibility through two layers:

\begin{enumerate}
  \item Create the role in Casdoor (\texttt{ROLE\_PLAY\_<GAME>} for game
        access, \texttt{ROLE\_ADMIN} for admin). The
        \snapfile{scripts/seed-roles.sh} script creates the per-game
        \texttt{ROLE\_PLAY\_*} roles idempotently; run it after
        deploying a new game.
  \item Insert a \texttt{game\_access} row with
        \texttt{(scope='role', value='ROLE\_PLAY\_<GAME>')} via the
        admin UI's \emph{Per-game access} card or by direct SQL.
\end{enumerate}

The admin dashboard's \emph{Grant role} dropdown reads the live Casdoor
catalog through \texttt{GET /api/admin/roles}; new roles appear as soon
as Casdoor's catalog has them.

\section{Adding a new SSE channel}

\begin{enumerate}
  \item Add a \texttt{Sinks.Many} field to \snapfile{GameStateService}
        with \texttt{autoCancel=false} (see the warning in
        chapter~\ref{ch:datamodel}).
  \item Add a \texttt{ServerSentEvent}-builder method that wraps the
        sink as a typed \texttt{Flux} with an \texttt{event:} field.
  \item Merge it into the \texttt{Flux.merge(...)} call inside
        \snapfile{GameStreamController} alongside the existing four
        channels and the keepalive.
  \item Register a per-type listener on the frontend by calling
        \texttt{eventSource.addEventListener('<event-name>', cb)} from
        \snapfile{GameStateService} or wherever the new channel is
        consumed.
\end{enumerate}

\section{Adding a new database column}

\begin{enumerate}
  \item Write a Liquibase changeset under
        \snapfile{backend/src/main/resources/db/changelog/} with the
        next sequential number. Postgres-only — Cassandra schemas live
        in \snapfile{infra/cassandra/cassandra-init.cql}.
  \item Include the changeset in \snapfile{db.changelog-master.xml}.
  \item Add the field to the matching R2DBC entity. The R2DBC mapping
        is reflective; \snapfile{R2dbcConvertersConfig} only needs new
        entries when adding a new entity, not a new column.
  \item For GraalVM native-image builds, run \texttt{make
        build-backend-native} and verify the reflection hints still
        resolve. If a new entity, regenerate the hints with the
        \texttt{native-image-agent} per the
        \reposrc{README.md} \emph{GraalVM native-image build} section.
\end{enumerate}

\chapter{Schema diagrams (appendix)}\label{ch:schema-diagrams}

The five auto-generated database diagrams below cover every table in
the platform. Each diagram is included on its own page sized to match
the source PDF's natural dimensions so a reviewer can pinch-zoom in
any PDF viewer without the content being rasterised down to thumbnail
size. Regenerated from the live Liquibase + Cassandra sources on every
\texttt{make docs-schema-diagram} run.

\begin{itemize}
  \item \textbf{BFF Postgres} --- catalog, identity, sessions, uploads,
        federation, role grants, friend graph, plus the per-game
        scope tables.
  \item \textbf{RPG-builder Postgres} --- Mudlet-shaped editor model
        (projects, areas, rooms, exits, items, NPCs, scripts) +
        per-job build state.
  \item \textbf{VN-builder Postgres} --- godot-2d-visual-novel-shaped
        editor model (projects, acts, characters, backgrounds,
        scenes, dialogue lines, branching choices) + per-job build
        state.
  \item \textbf{Cassandra} --- the cache, audit log, and rate-limit
        event log (active\_sessions, controller\_lease,
        character\_position\_cache, chapter\_membership,
        session\_membership, session\_view\_state, rate\_limit\_events,
        aggregate\_counters, pod\_heartbeats).
  \item \textbf{Combined view} --- every table across all four stores,
        with the cross-cluster foreign keys drawn explicitly.
\end{itemize}

% The script emits one \afterpage hook + \includepdf pair per diagram.
% \afterpage queues actions FIFO, so each shipout picks up the right
% \pdfpagewidth/\pdfpageheight for its diagram. The whole section
% lives in this appendix so the over-sized pages don't disrupt earlier
% chapters' typography.
\input{../diagrams-rendered/build/schema-diagram-includes-arch.tex}


% ── Back matter ───────────────────────────────────────────────────────────
\backmatter

\chapter{Where to next}

This reference covers the system as it exists at the head of the
\reposrc{feature/architecture-publishing} branch. For deeper dives:

\begin{itemize}
  \item \reposrc{README.md} — the canonical setup and operations guide.
  \item \reposrc{docs/PLAN-module-federation-and-native.md} — the
        eight-section plan that produced the federated topology.
  \item \reposrc{docs/GODOT_DEVELOPER_GUIDE.md} — game-developer-facing
        contract for slugs, presets, and uploads.
  \item Generated Javadoc HTML at
        \reposrc{docs/publishing/javadoc-output/index.html} after running
        \snapcommand{make docs-javadocs}. The same HTML is published as a
        GitHub Pages artefact when CI runs.
  \item PowerPoint walkthroughs under
        \reposrc{docs/publishing/slides/build/} after running
        \snapcommand{make docs-slides}.
\end{itemize}

\noindent
The shortest path back to live code is
\href{https://github.com/hugmesh-net/hugmesh-games}{github.com/hugmesh-net/hugmesh-games}.

\end{document}
